\chapter*{Liste des abréviations}
\label{chap:abbreviations}
\addstarredchapter{Liste des abréviations}
\markboth{Liste des abréviations}{}
\begin{itemize}
\item[•] \textbf{AVX2 :} \emph{Advanced Vector Extensions 2} : instructions des processeurs qui traitent 256 bits (32 octets) à la fois
\item[•] \textbf{CT :} \emph{Constant Time} (temps constant)
\item[•] \textbf{DSE :} \emph{Dead Store Elimination} (passe de GCC)
\item[•] \textbf{FRE :} \emph{Full Redundancy Elimination} (passe de GCC)
\item[•] \textbf{GCC :} \emph{GNU Compiler Collection}
\item[•] \textbf{GFNI :} \emph{Galois Field New Instructions}
\item[•] \textbf{IC95 :} intervalle de confiance à 95\,\%
\item[•] \textbf{KAT :} \emph{Known-Answer Test} : test de « réponses connues » (on compare à des valeurs officielles)
\item[•] \textbf{LIM :} \emph{Loop Invariant Motion} (passe de GCC)
\item[•] \textbf{MAYO :} schéma de signature multivariée de type Oil-and-Vinegar
\item[•] \textbf{FIPS :} \emph{Federal Information Processing Standards} : normes fédérales américaines de traitement de l'information
\item[•] \textbf{MQ :} \emph{Multivariate Quadratic} : équations de degré 2 à plusieurs variables
\item[•] \textbf{NIST :} \emph{National Institute of Standards and Technology} : Institut national américain des normes et de la technologie
\item[•] \textbf{OV, UOV :} \emph{(Unbalanced) Oil and Vinegar}
\item[•] \textbf{PRE :} \emph{Partial Redundancy Elimination} (passe de GCC)
\item[•] \textbf{SIMD :} \emph{Single Instruction, Multiple Data}
\item[•] \textbf{SRA :} \emph{Scalar Replacement of Aggregates} (passe de GCC)
\item[•] \textbf{TSC :} \emph{Time Stamp Counter} : compteur de ticks du processeur
\end{itemize}
\clearpage
